A central methodological contribution of this paper is that, by matching every listing against the ADEME registry, we can directly observe whether a seller is informed ($K=1$, audit exists) or uninformed ($K=0$, no audit on record). Crossing this latent information status with the seller's disclosure choice $d \in \{0,1\}$ yields a four-cell classification of which only three cells are economically distinct: truthful disclosers ($K=1, d=1$), strategically hiding sellers ($K=1, d=0$), and genuinely uninformed sellers ($K=0, d=0$).\footnote{The fourth cell ($K=0$, $d=1$) comprises sellers who display an EPC score for which no matching record appears in the public ADEME registry. These listings are treated as truthful disclosers, consistent with the no-falsification assumption: a score published by the owner is taken at face value, and only omission is treated as a strategic act. Two institutional features explain the size of this group. First, the public registry contains only certificates issued from July 2021 onward, whereas a certificate issued earlier could still be legally valid at the time of listing. Under \emph{d\'ecret} n\textsuperscript{o}~2020-1610 of 17 December 2020, certificates issued between 1 January 2013 and 31 December 2017 remained valid until 31 December 2022, and those issued between 1 January 2018 and 30 June 2021 until 31 December 2024, while certificates issued from 1 July 2021 are valid for ten years. A seller listing during our sample window could therefore display a still-valid certificate that is absent from the registry by construction. Second, the model assumes no falsification. Consistent with the first explanation, the share of listings in this cell falls monotonically as the registry accumulates post-reform certificates: from $62.5\%$ of apartment listings in 2021Q3 to $18.4\%$ in 2024Q4, and from $73.1\%$ to $43.7\%$ for houses.} Figure~\ref{fig:fact3_decomposition} reports the resulting decomposition for both segments. In apartments (panel~(a)), $79.8\%$ of listings are truthful disclosers, $16.7\%$ are genuinely uninformed, and $3.5\%$ are strategically hiding. In houses (panel~(b)), the corresponding shares are $90.4\%$, $8.5\%$, and $1.1\%$. Houses systematically display higher baseline disclosure and a smaller strategic-hiding share, but the qualitative structure is identical: a non-trivial mass of strategic hiders coexists with a much larger mass of genuinely uninformed sellers in both segments.
